\section{Fixed-factor conic strata in degrees five and six}
\label{app:degree-five-six}

This appendix pushes the conic leading-image method beyond degree four.  It
does not classify degrees five or six: the theorem isolates an open
basepoint boundary and a resonant primitive sextic case.

Let the leading homogeneous map have the squarefree fixed-factor form
\[
\begin{aligned}
H_D&=G(A^2,AB,B^2),\\
\deg G&=g>0,\qquad \deg A=\deg B=e,\qquad D=g+2e.
\end{aligned}
\]
Assume that \(A,B\) are a regular sequence and that some irreducible
component of \(G=0\) sees a nonconstant ratio \(A/B\).  These are the
weaker invariant-gap hypotheses used in the next lemma; they do not imply
that the base locus of \((A,B)\) is disjoint from \(G=0\).

\begin{lemma}[Invariant-degree gap]
\label{lem:fixed-factor-invariant-gap}
For
\[
\Omega=A\,dG\wedge dB-B\,dG\wedge dA+2G\,dA\wedge dB,
\]
let \(\delta\) be the associated homogeneous derivation.  Every nonzero
homogeneous \(P\) satisfying \(\delta(P)=0\) has degree divisible by \(D\).
\end{lemma}

\begin{proof}
Choose an irreducible factor \(\Gamma\mid G\), of degree \(h\), on which
\(A/B\) is nonconstant, and write \(P=\Gamma^rQ\) with
\(\Gamma\nmid Q\).  On the function field of \(\Gamma=0\), both the induced
derivation and the Euler derivation annihilate \(A/B\); hence
\(\bar\delta=cE\) for a nonzero rational function \(c\).  From
\(\delta(GB^2)=0\) one gets
\[
\overline{\delta\Gamma/\Gamma}=-c(D-h).
\]
Reducing \(\delta(\Gamma^rQ)=0\) modulo \(\Gamma\) now yields
\[
c(\deg P-rD)\bar Q=0.
\]
Thus \(\deg P=rD\).
\end{proof}

Write a Keller jet as
\[
K_\varepsilon
=H_D+\varepsilon H_{D-1}+\cdots+\varepsilon^{D-2}H_2
+\varepsilon^{D-1}LX
\]
and set
\[
\Phi(Y)=Y_1Y_3-Y_2^2,\qquad
\Phi(K_\varepsilon)=\sum_{j\ge1}\varepsilon^jS_j.
\]
The chain-rule identity
\[
\operatorname{Jac}(K_{\varepsilon,1},K_{\varepsilon,2},\Phi(K_\varepsilon))
=K_{\varepsilon,1}\det JK_\varepsilon
\]
and \cref{lem:fixed-factor-invariant-gap} imply inductively that
\[
S_1=\cdots=S_{D-1}=0.
\]
Indeed, \(\deg S_j=2D-j\), which lies strictly between \(D\) and \(2D\).

\begin{theorem}[Basepoint-free fixed-factor conic exclusion]
\label{thm:degree-five-six-fixed-factor}
Under the preceding hypotheses, assume in addition that
\[
V_{\mathbf P^2}(G,A,B)=\varnothing.
\]
Then no nonautomorphic Keller map has one of the following leading-form
types:
\[
\begin{array}{c|c}
D=5&(g,e)=(1,2),(3,1)\\
D=6&(g,e)=(2,2),(4,1).
\end{array}
\]
\end{theorem}

\begin{proof}[Proof idea]
The additional projective basepoint-free hypothesis is used here: it says
that \(A,B\) form a regular sequence in the homogeneous coordinate ring
\(\C[x,y,z]/(G)\).  The vanishings \(S_1,\ldots,S_{D-1}\) combine with the
Koszul syzygies of \((A,B)\) and the Hilbert--Burch syzygies of
\((A,B)^2\).  They lift the
deformed rank-one symmetric matrix to
\[
\widehat{\mathsf M}_\varepsilon
=a_\varepsilon P_\varepsilon P_\varepsilon^T
+\varepsilon^3\operatorname{sym}(P_\varepsilon,Q_\varepsilon)
\]
after a source-independent target correction.  Consequently
\[
\Phi(\widehat K_\varepsilon)
=-\varepsilon^6(P_\varepsilon\wedge Q_\varepsilon)^2.
\]
The chain-rule identity forces every non-\(\varepsilon\) irreducible factor
of \(P_\varepsilon\wedge Q_\varepsilon\) to divide both of the first two
coordinates and then the Keller determinant.  This is impossible.  The
remaining possibility is a pure power of \(\varepsilon\); comparison of the
lowest coefficients makes the wedge vanish and again forces the linear
determinant to be zero.

Under \texttt{code/degree-five-six/}, the file
\path{conic_fixed_factor_d5_d6_proof_note.md} writes out the four
graded lift identities.  The script
\path{verify_conic_hensel_d5_d6.py} replays them exactly.
\end{proof}

The invariant-degree gap itself remains valid under the weaker hypotheses
stated before \cref{lem:fixed-factor-invariant-gap}.  What happens when
\(V_{\mathbf P^2}(G,A,B)\ne\varnothing\) is not settled by the lifting
argument above and remains an open boundary problem.

The \((g,e)=(3,1)\) quintic case admits a stronger boundary analysis.

\begin{theorem}[Three-variable cubic-factor quintics]
\label{thm:quintic-three-variable-cubic-factor}
Let \(G\) be a squarefree cubic genuinely depending on \(x,y,z\).  There is
no Keller jet with leading form
\[
G(x,y,z)(x^2,xy,y^2).
\]
Thus the only unresolved cubic-factor overlap of this shape is the binary
case \(G\in\C[x,y]\).
\end{theorem}

\begin{proof}[Proof structure]
At the base point \([0:0:1]\), the four conic normal equations split
according to the multiplicity of \(G=0\) and whether the first normal linear
form vanishes there.  The unit branch is eliminated by the rank of the
tangent cone and, in the rank-one case, by two incompatible coefficients.
In the regular branch a source translation removes the first syzygy.
The remaining cusp, node, transverse line--conic, and tangent line--conic
normal forms each collapse under two or three displayed coefficients of the
Keller determinant.  Multiplicity three forces \(G\) to be binary.

The complete normal-form ledger and decisive coefficients are recorded in
\path{quintic_cubic_factor_basepoints.md}, under
\texttt{code/degree-five-six/}.  The companion exact script replays every
coefficient identity used in the case table.
\end{proof}

\begin{question}[Binary cubic-factor overlap]
\label{q:binary-quintic-overlap}
Can a quintic Keller map have leading form
\[
G(x,y)(x^2,xy,y^2)
\]
for a squarefree binary cubic \(G\)?
\end{question}

\begin{question}[Fixed-factor basepoints in degree five]
\label{q:degree-five-fixed-factor-basepoints}
Can the remaining \((g,e)=(1,2)\) quintic branches, where the base locus of
\((A,B)\) meets \(G=0\), satisfy the Keller equations?
\end{question}

\begin{question}[Primitive sextic conic stratum]
\label{q:primitive-sextic-conic}
Can a degree-six Keller map have primitive conic leading form
\[
(A_3^2,A_3B_3,B_3^2)
\]
with no nonconstant fixed factor?  Here degree three is resonant, so the
invariant-gap argument does not apply.
\end{question}
